\section{The filtered projection for a down position followed by an up position}

% BEGIN P82U1
\paragraph{P82U1. Testing an up vector on the two physical lattices.}
Take the backward positions $d<u$ of A82N4, with $d$ down and
$u$ up, and the actual onto map of tested lattices. Keep the same
DVR $O$ throughout. Write $D=P_d^{b_d}$, $P=P_u^{b_u}$,
$A=O[[v]]$, $\Theta=\prod_{s<\ell}P_s^g$, $a=2\sigma+rr$,
and $f_0$ for the map in the constant frames. For any
$e\in U_{2,u}$ the target numerator
\[
                    y_2=(\Lambda/P)D e
\]
lies in $\mathrm{NN}_2$: $y_2/\Lambda$ is regular away from
the up position, vanishes to order $b_d$ at the down position,
and its only possible pole at the up position has coefficients
in $U_{2,u}$. We claim that there is an integral numerator
$y_1\in A\otimes T_1$ with
\[
 \Theta^2 y_1\in\mathrm{NN}_1,\qquad
   f_0 y_1-p^a y_2\in \Lambda(A\otimes T_2).
\]
Here $T_i$ denotes the constant coordinate lattice of A82N1.

To prove this, apply J82F2 to $y_2$. Its physical partner
$w_2\in W_2$ agrees with $p^\sigma y_2/\Lambda$ in period
coordinates up to the error of order $d_s-H_s$ at each root
in level $s$, where $d_s$ is defined in J82F1. At this level
write $N_{i,s}$ for the true Taylor lattice and $B_{i,s}$ for
the rounded one of A82N1. D82P3 gives
$P_s^gB_{i,s}\subset N_{i,s}$ and
$P_s^gN_{i,s}\subset B_{i,s}$, and the error belongs to
both lattices. These facts concern the actual representation:
the two lattices have the same adapted flag, with its true jumps
and with those jumps rounded to $0$ or $\pm b_s$, respectively.
Their differences in exponent are at most $g$. The filtration
at this level is transported from $c-s$ by $\varphi^s$.
Consequently $\Theta w_2$ is in every true Taylor lattice.
D82P1 gives $\Theta w_2\in N_{2,c}$ integrally, since a
coordinate denominator in $W_2$ is a product of the distinguished
level factors and Taylor regularity at every root cancels it.

Use R82W4 on this vector at the terminal cut $c$, then divide its
source lift by $\Theta$. This yields $w_1\in\Theta^{-1}N_{1,c}$
with $f(w_1)=p^{rr}w_2$. This is an element of $W_1$ by D82P2,
since $\Lambda/\Theta$ is the product of the successive upper
Frobenius height factors. J82F2 in the other direction gives
an integral $y_1$ whose projection $y_1/\Lambda$ differs from
$p^\sigma w_1$ only by the specified errors. At each level
$\Theta w_1$ is in the true lattice, hence
$\Theta^2 w_1$ is in the rounded lattice. The errors are in
the rounded lattice as well, so A82N1 gives the first claim.
By functoriality J82F3, the differences composing
$(f_0 y_1-p^a y_2)/\Lambda$ are in the order $d_s-H_s$ error
space. This order is nonnegative: on the terminal range all
positions are thin with $H_s\le2g$ and $d_s=3g$, and on the
other range $d_s=(p-1)M>h+2g\ge H_s$. Thus the quotient is
regular at every root. Weierstrass division by the level factors
as in A82N1 gives the asserted integral divisibility by $\Lambda$.
% END P82U1

% BEGIN P82U2
\paragraph{P82U2. The valuation bound in the two-level quotient.}
Suppose that in $A$ one has
\[
       p^a\Theta^2 P_d^{b_d}\in (\lambda^k,P_u^{b_u})
\]
where $\lambda$ is a uniformizer. Put $e_O=v_\lambda(p)$
and $A_0=a+2g\ell$. The guards of I82B2--I82B3 give
$A_0\le a_0$, $2a_0<b_u$, and $2b_d\le p b_u$. Then
\[
             k/e_O\le \frac{A_0 b_u}{b_u-b_d/p}\le 2a_0 .
\]
Indeed, give $v$ Gauss weight $t/p^u$ with $0<t<1$ and
normalize $v_p(p)=1$. For an integral series the weighted
valuation is the minimum of its coefficient valuations plus
degree times this weight. This minimum is attained in finite
degree, and the least degrees attaining the minima in two
factors give the least attaining degree in their product.
Thus the valuation is multiplicative. Each distinguished
$P_s$ is a unit times a monic polynomial of degree $p^s$ with
lower coefficients divisible by $p$ and constant term of
valuation one. Indeed it has coefficients in $\mathbb Z_p$,
residue order $p^s$ and constant term $p$ by A82N1, so its
Weierstrass polynomial over $\mathbb Z_p$ has these properties.
Hence its weighted valuation is $\min(1,p^{s-u}t)$.

It follows that the left side in the ideal inclusion has valuation
at most $A_0+(b_d/p)t$, since $d<u$. Every element of the ideal
has valuation at least $\min(k/e_O,b_u t)$. Now take
$t>A_0/(b_u-b_d/p)$, still less than one; this interval exists
by the two strict or weak size guards above. The upper bound is
strictly less than $b_u t$. Thus $k/e_O\le A_0+(b_d/p)t$.
Let $t$ decrease to the endpoint. This proves the claimed
inequality, and also rules out inclusion in $(P_u^{b_u})$ alone.
% END P82U2

% BEGIN P82U3
\paragraph{P82U3. The native filtered lift for the last sign order.}
The true map on crystalline components preserves filtration, so
$f_0 U_{1,u}\subset U_{2,u}$ by J82F3 and the transported flag
in P82U1. Take Smith coordinates for this inclusion inside the
free lattice $U_{2,u}$. Since $U_{2,u}$ is saturated in $T_2$,
each Smith coordinate functional extends to $T_2$. Fix one, say
$l$, take its basis vector $e\in U_{2,u}$ with $l(e)=1$,
and let $I=l(f_0 U_{1,u})$, an ideal of $O$.

Use P82U1 for this $e$ and form $q=\Theta^2 y_1/(\Lambda/P)$.
This is integral by the first assertion there: for the
homogeneous model all levels except $u$ require regularity,
and the only allowed pole at $u$ has order at most $b_u$.
Its component modulo $U_{1,u}$ is divisible by $P$, again by
A82N1 and saturation. Thus $q\in A U_{1,u}+P(A\otimes T_1)$.
Multiplying P82U1's second assertion by $\Theta^2/(\Lambda/P)$
gives
\[
            f_0 q\equiv p^a\Theta^2D e\pmod{P(A\otimes T_2)} .
\]
Application of $l$ yields $p^a\Theta^2 D\in IA+PA$.
By P82U2, $I\ne0$, and when $I=(\lambda^k)$ it contains
$p^{2a_0}$. This argument works in every Smith coordinate, so
\[
                p^{2a_0}U_{2,u}\subset f_0(U_{1,u}).
\]
It is precisely $\mathrm{CLift}$ on the actual tested filtered
coordinate lattices. A82N5 now gives $\mathrm{ConstLift}$ on
this down-before-up arc too. Consequently the normalized integral
constant numerators lift at the stated power for every sign
order of a marked row B arc.
% END P82U3
